\documentclass[11pt]{article}
\usepackage[margin=1in]{geometry}
\usepackage{enumitem}
\usepackage[T1]{fontenc}
\setlength{\parskip}{0.5em}
\setlength{\parindent}{0pt}
\begin{document}

\hfill Andrew H. Bond\\
\hfill Department of Computer Engineering\\
\hfill San Jos\'e State University\\
\hfill \texttt{andrew.bond@sjsu.edu} \textbullet\ ORCID 0009-0003-2599-6158

\bigskip
To the Editors, \emph{SIAM Journal on Mathematics of Data Science}

\bigskip
Dear Editors,

Please consider the enclosed manuscript, ``Keep the Angle: A Geometry-Preserving
Basis in Spectral Embeddings,'' for publication in \emph{SIMODS}.

\textbf{Summary.} The commute-time (resistance) embedding built from the low
eigenvectors of a graph Laplacian degenerates on large graphs (von~Luxburg--Radl--Hein).
We show the geometry does not vanish but is carried by the \emph{angular}
coordinate---the row-normalization of Ng--Jordan--Weiss spectral clustering,
here given a metric-geometric rather than a cluster-separation reading. The
contribution is organized as four claims: (1) an exact normalization identity
locating the mechanism in \emph{tangential noncollapse}, and a deterministic
angular-transfer theorem; (2) a conditional graph-to-continuum transfer, carried
by the random-walk eigenvectors and the sup-norm eigenvector rates of
Calder--Garc\'ia Trillos--Lewicka, together with an \emph{unconditional}
bi-Lipschitz statement on the flat torus; (3) a filter dichotomy---uniform-in-truncation
for the heat filter, scale-dependent for the commute filter, with an exact
rank-collapse theorem for the plain filter on the circle; and (4) a controlled
empirical study across five substrate families. We believe the topic---when and
why row-normalized spectral embeddings recover intrinsic geodesic order---is
squarely within the journal's scope.

\textbf{Originality and exclusivity.} This manuscript is original, has not been
published, and is not under review at any other journal or conference. We intend
to post it as a preprint, consistent with SIAM policy.

\textbf{Related work by the author (overlap disclosure).} In the interest of full
transparency we note two of the author's own companion artifacts that the
manuscript cites for context, and which are \emph{not} overlapping submissions:
\begin{itemize}[nosep,leftmargin=1.4em]
  \item \emph{Observation Theory: geometry, distortion, and reliability as
  properties of observation} (technical report / preprint, Zenodo
  \texttt{10.5281/zenodo.21423315}). This is a broader, consumer-relative
  framework; the present paper uses only its ``tangential-noncollapse (A2)''
  vocabulary as an organizing label and shares no theorem with it.
  \item \emph{TurboQuant-Pro} (open-source software and accompanying technical
  note). Cited once, as a concrete boundary example (per-vector angular
  quantization of attention \emph{keys} fails because scale is signal there).
  It is a software project, not a journal submission.
\end{itemize}
Neither is submitted or accepted elsewhere; copies can be provided to the editors
and referees on request.

\textbf{Revision note.} The manuscript underwent detailed pre-submission review;
the principal changes are summarized in the point-by-point response below, and the
two new results (the $S^1$ rank-collapse theorem and the random-walk transfer
repair) were subjected to an independent cold re-derivation before submission.

Thank you for your consideration.

\bigskip
Sincerely,\\
Andrew H. Bond

\newpage
\begin{center}\large\textbf{Response to review}\end{center}

We are grateful for the careful reading. Below we address each major point; the
verdict was \emph{major revision}, and every item is now resolved. Section and
label references are to the revised manuscript.

\textbf{1. Omitted prior use of the normalization.} We now cite Sharma--Horaud--Mateus
(2012), who apply unit-hypersphere normalization directly to the commute-time
embedding for shape registration, and we reframe the novelty accordingly: the
normalization is not new; the metric-geometric account of \emph{when its angular
distances recover intrinsic geodesic order} is our contribution (\S1).

\textbf{2. Proposition 14 did not prove rank collapse.} Correct---distance
concentration does not imply rank collapse. We split the result into an honest
general-manifold \emph{concentration} proposition and a new \emph{exact} theorem
(Theorem, \texttt{thm:s1collapse}) proving that on $S^1$ both Spearman and Kendall
correlation with geodesic distance vanish, via the Dirichlet kernel and two-scale
equidistribution. This is materially stronger than weakening the claim.

\textbf{3. Theorem 8 normalization gap ($L_{\mathrm{sym}}$ vs.\ $L_{\mathrm{rw}}$).}
Adopted the suggested repair. The angular map is invariant to the exact per-node
factor $\sqrt{d_i}$ ($\Psi_i^{\mathrm{sym}}=\sqrt{d_i}\,\Psi_i^{\mathrm{rw}}$), so
we carry the sup-norm transfer on the random-walk eigenvectors---for which the
cited rates are stated---and read off the algorithm's $L_{\mathrm{sym}}$ angle by
exact cancellation. No degree-concentration hypothesis is needed; the theorem now
holds under non-uniform sampling. Step~3 is written as an external lemma
(\texttt{lem:transfer}) composing named rate theorems (Calder--Garc\'ia Trillos
for eigenvalue/gap stability; Calder--Garc\'ia Trillos--Lewicka for the $L^\infty$
and $C^{0,1}$ eigenvector rates), making the theorem unconditional in $m$ at each
fixed $m$; the $m$-dependence of those constants is stated explicitly as the
bridge to the open problem.

\textbf{4. The central conjecture was quantified too broadly.} Conjecture is now
scoped to an admissible mode schedule (multiplet-closed $m$ at a spectral gap,
$m\ge m_0(M)$, the spectral-window condition for growing $m$, pairs above the
resolution scale), with ``uniform'' meaning a bound over that family. We name the
hypothesis \texttt{(H1)-uniform}, prove its two endpoints (unconditional on the
torus; a continuum theorem for the heat filter, with the commute filter shown
\emph{not} m-uniform at fixed scale), and state the single sharp open problem---a
gap-uniform lower bound on the smallest singular value of the truncated-eigenmap
Jacobian (\texttt{conj:sigmamin})---whose resolution closes the metric form.

\textbf{5. Sliding between two embeddings; $d\ge3$ vs $d\ge2$.} We now state the
radial claim at three fixed levels and never mix them: proved for the full
unnormalized embedding at $d\ge3$; empirical for the truncated normalized radius
the algorithm computes ($\rho\approx0.79$--$0.87$ vs.\ $0.98$ for the full
radius); borderline at $d=2$. A ``Status of the radial claim'' block records that
the gap between the proved and empirical levels is exactly the truncation, which
no theorem here bridges.

\textbf{6. Proposition 13 remainder.} We supply the required differentiated
pointwise Weyl law as a displayed lemma (\texttt{lem:dweyl}), citing H\"ormander
for the on-diagonal projector and Canzani--Hanin for the mixed-derivative
remainder (both already in the bibliography), and add a cutoff-insensitivity
remark noting the $S^1$ case needs no Weyl input.

\textbf{7. The ``scaling law'' is not a law.} Retitled to ``an exploratory
dimension--fidelity trend.'' The five rule-level summaries are now the primary
table; the twenty-point fit is presented as a descriptive trend line with its
rule-clustered CI, and Proposition (rank) is repositioned as \emph{consistent
with, and not deriving,} the trend.

\textbf{8. Methods too compressed; graph vs.\ manifold; replication.} Added a
per-section protocol table; we report both the measured dimension $d_G$ and the
ground-truth $d_M$ where it exists, noting the finite-sample estimator bias shifts
the axis but not the ordering; and the headline uncertainty is reported over
independent graph draws (pair-resampling intervals are labeled as optimistic
floors).

\textbf{Additional comments.} Observer-theory and Wolfram framing are removed from
the title and abstract and confined to a single discussion paragraph; the abstract
is rewritten to $\le250$ words with references spelled out; terminology is
softened from ``basis'' to ``representation'' where literally apt; the figure
caption/text discrepancies are reconciled and min--max ranges are labeled as
ranges rather than confidence intervals.

\end{document}
